\documentclass[10pt,a4paper]{amsart}
\usepackage[margin=19mm]{geometry}
\usepackage{amsmath,amssymb,mathtools}
\usepackage[scr=rsfs,cal=euler]{mathalfa}
\usepackage{xfrac}
\usepackage[unicode,hidelinks,bookmarks=false]{hyperref}
\newcommand{\ie}{i.e.\ }
\newcommand{\cf}{cf.\ }
\newcommand{\resp}{respectively\ }
\newcommand{\Subsection}{\subsection}
\providecommand{\nobreakdash}{\nobreak\hbox{-}}
\setlength{\emergencystretch}{2em}
\multlinegap0pt
\usepackage{bm}
\usepackage{bbm}
\usepackage{dsfont}
\usepackage{xspace}
\usepackage{cancel}
\usepackage{mathtools}
\usepackage{amsthm}
\makeatletter
\@ifundefined{theorem}{\newtheorem{theorem}{Theorem}}{}
\@ifundefined{lemma}{\newtheorem{lemma}[theorem]{Lemma}}{}
\makeatother
\newcommand{\Z}{\ensuremath{\mathbb{Z}}}
\newcommand{\vr}[2]{\mathrm{VR}(#1;#2)}
\title[Vietoris-Rips complexes of torus grids]{Vietoris-Rips complexes of torus grids}
\author{Henry Adams, Adenike Yeside Adetowubo, Hector Barriga-Acosta, Ziqin Feng, John Sterling}
\date{Source: arXiv:2502.07134v2 (CC BY 4.0)}
\begin{document}
\maketitle

\noindent\emph{Source and licence.} Vietoris-Rips complexes of torus grids. Version arXiv:2502.07134v2 (CC BY 4.0), distributed under the Creative Commons Attribution 4.0 International licence, as recorded in the arXiv licence metadata for this exact version and shown on the abstract page https://arxiv.org/abs/2502.07134v2. Full rights notice: licenses/arxiv-2502.07134-CC-BY-4.0.txt.\par\smallskip

\noindent\textit{Editorial scope.} Counted unit: the Lemma whose source label is \texttt{lem:facets\_3kor3k-1} in the article (source0.tex lines 1689--1693, source label \texttt{lem:facets\_3kor3k-1}), delivered together with the article's own complete original proof of it (source0.tex lines 1695--1741, one \texttt{proof} environment of 47 lines). No mathematics is written, repaired or re-derived: statement and proof are the article's own text line for line. Required context retained uncounted: lemma:imageoffacets (lines 1047--1051); lemma:prefacets (lines 1065--1070); lemma:emptyintersection (lines 1109--1113). Every definition, display and label of those lines is retained unchanged. Omitted from this package and not redistributed: 1--1046, 1052--1064, 1071--1108, 1114--1688, 1694--1694, 1742--2655. Nothing inside a retained excerpt is omitted or altered: every retained body is delivered exactly as the article's lines are written, so no omission receipt of any kind accompanies this package. Citations: the retained excerpts carry no citation command at all, so no bibliography entry of the article is transcribed and no reference is redistributed. Reference adaptation: none was needed and none was made; every reference inside the delivered unit resolves inside it, so no unresolved reference or citation is produced. Printed slips kept literal: no printed word, symbol, hypothesis or number of the counted statement or of its proof was altered. Layout and class: the article's own class is \texttt{amsart} with packages including inputenc, amsmath, amsthm, amsfonts, amssymb, graphicx, caption, subcaption, xcolor, chngpage, color, tikz, hyperref, fullpage; the wrapper is the lane's pinned \texttt{amsart} editorial wrapper at 10pt on A4, which restates the theorem-like environments the retained text needs and adds the article's own macro definitions that the retained text needs (\texttt{\textbackslash Z}, \texttt{\textbackslash vr}), with extra wrapper packages (bm, bbm, dsfont, xspace, cancel, mathtools). The delivered numbering is the unit's own. Assets: no retained span contains a figure environment, a graphics-inclusion command, a PDF-page inclusion, a table or a TikZ picture; no image file is copied into this package, the package declares no asset dependency and no image file is redistributed with it. Licence: version arXiv:2502.07134v2 of this article is distributed under the Creative Commons Attribution 4.0 International licence, as recorded in the arXiv licence metadata for this exact version and shown on the abstract page https://arxiv.org/abs/2502.07134v2. A full rights notice is delivered at licenses/arxiv-2502.07134-CC-BY-4.0.txt. Characters: every retained excerpt is pure ASCII, so the delivered engine, XeLaTeX with its default Latin Modern text font, reports no missing character.
\begin{lemma}
\label{lemma:imageoffacets}
Let $\sigma$ be a facet in the complex $\vr{\Z^2}{k}$.
For $n>2k+1$, $\pi_n(\sigma)$ is a facet in $\vr{\pi_n(\Z^2)}{k}$.
\end{lemma}

\begin{lemma}
\label{lemma:prefacets}
Suppose that $n\geq 3k-1$ and $k\geq 1$.
Let $x,y\in\Z^2$ with $x\neq y$, with $3k-n< d(x,y)\leq k$, and with $\pi_n(B_{\Z^2}[x, k]\cap B_{\Z^2}[y, k]) =\pi_n(B_{\Z^2}[x, k])\cap \pi_n(B_{\Z^2}[y, k])$.
Then any facet $\tau$ in $\vr{\pi_n(\Z^2)}{k}$ containing $[x]$ and $[y]$ is of the form $\pi_n(\sigma)$ for some facet $\sigma\in \vr{\Z^2}{k}$.
\end{lemma}

\begin{lemma}
\label{lemma:emptyintersection} 
Fix different $x=(x_1,x_2)$ and $y=(y_1,y_2)$ in $\Z^2$.
If $|x_1-y_1|<n-2k$ and $|x_2-y_2|<n-2k$, then $\pi_n (B_{\Z^2}[x,k]\cap B_{\Z^2}[y,k]) =\pi_n(B_{\Z^2}[x,k]) \cap \pi_n(B_{\Z^2}[y,k])$.
\end{lemma}

\begin{lemma}
\label{lem:facets_3kor3k-1}
For any $k\geq 2$, the collection of facets in $\vr{T_{3k,3k}}{k}$ is $M_{3k,k}\cup N_{3k,k}$.
Also, for any $k\geq 3$, the collection of facets in $\vr{T_{3k-1,3k-1}}{k}$ is $M_{3k-1,k}\cup N_{3k-1,k}$.
\end{lemma}

\begin{proof}
We have $M_{3k,k}\subseteq M(\vr{T_{3k,3k}}{k})$ by Lemma~\ref{lemma:imageoffacets} (with $3k=n>2k+1$ since $k\ge 2$).
Also, we have the containment $N_{3k,k}\subseteq M(\vr{T_{3k,3k}}{k})$.
Similarly, we have $M_{3k-1,k}\subseteq M(\vr{T_{3k-1,3k-1}}{k})$ by Lemma~\ref{lemma:imageoffacets} (with $3k-1=n>2k+1$ since $k\ge 3$), along with the containment $N_{3k-1,k}\subseteq M(\vr{T_{3k-1,3k-1}}{k})$.

Next we pick a facet $\tau$ in the complex $\vr{T_{3k,3k}}{k}$.
We'll show that $\tau$ is in $M_{3k,k}\cup N_{3k,k}$.

First we suppose that there exist $[x] =([a], [b])$ and $[y]=([c], [d])$ in $\tau$ with $[a]\neq [c]$ and $[b]\neq [d]$.
Without loss of generality, we assume that $x=(a, b)$ and $y=(c, d)$ and $d(x, y)=d([x], [y])\le k$.
Clearly $a\neq c$ and $b\neq d$; hence $1\leq d(a, c)\leq k-1$ and $1\leq d(b, d)\leq k-1$.
By Lemma~\ref{lemma:emptyintersection} (with $n-2k=k$), we have $\pi_n (B_{\Z^2}[x, k]\cap B_{\Z^2}[y, k]) =\pi_n(B_{\Z^2}[x, k]) \cap \pi_n(B_{\Z^2}[y, k])$.
Then by Lemma~\ref{lemma:prefacets}, we have $\tau=\pi_n(\sigma)$ for some facet $\sigma$ in $\vr{\Z^2}{k}$.
Therefore $\tau\in M_{3k,k}$.

%Then, we suppose that any pair of elements in $\tau$ share the same first coordinate or second coordinate.
%This implies that all elements in $\tau$ have the same first coordinate or second coordinate.
%We assume thatthe first coordinate of all elements in $\tau$ is $[a]$.

Otherwise, all elements in $\tau$ have the same first coordinate or second coordinate.
Without loss of generality we assume that the first coordinate of all elements in $\tau$ is $[a]$.
The full subcomplex of $\vr{T_{3k,3k}}{k}$ whose vertices have first coordinate $[a]$ is isomorphic to the clique complex $\vr{C_{3k}}{k}$, where $C_{3k}$ is the cyclic graph with $3k$ vertices.
Of the maximal simplices in this copy of $\vr{C_{3k}}{k}$, only the ones of the form $\{([a], [b]), ([a], [b+k]), ([a], [b+2k])\}$ for some $0\leq b\leq 3k-1$ are maximal in $\vr{T_{3k,3k}}{k}$.
Hence $\tau\in N_{3k,k}$.
This completes the proof of the first statement of the lemma. 

\medskip 

For the second statement, it remains to show that a facet $\tau$ in $\vr{T_{3k-1,3k-1}}{k}$ is in $M_{3k-1,k}\cup N_{3k-1,k}$.
Consider the case where there exist $[x] =([a], [b])$ and $[y]=([c], [d])$ in $\tau$ with $[a]\neq [c]$ and $[b]\neq [d]$.
Then the same argument as before applies, \emph{except} we cannot apply Lemma~\ref{lemma:emptyintersection} when $d(x, y)=k$ with either $d(a, c)=1$ or $d(b, d)=1$.
Without loss of generality, assume $d(a, c)=1$ and $d(b, d)=k-1$.
We show that still in this case $\pi_n (B_{\Z^2}[x, k]\cap B_{\Z^2}[y, k]) =\pi_n(B_{\Z^2}[x, k]) \cap \pi_n(B_{\Z^2}[y, k])$, in which case $\tau$ is in $M_{3k-1,k}$ by Lemma~\ref{lemma:prefacets}.
Moreover, we assume that $a<c$ and $b<d$, and then $c=a+1$ and $d=b+k-1$.
Other cases can be proved similarly.
Pick an arbitrary $[z]\in B_{\Z^2}[y, k]$ with $d((z_1, z_2), y)=|z_1-c|+|z_2-d|\leq k$. 
It is sufficient to show that $d(z', x)>k$ for any $z'\in [z]$ with $z'\neq (z_1, z_2)$.
We show this holds for $z'=(z_1+3k-1, z_2)$, and similar approach can be applied for $z'$ in different forms.
Notice that \[d(z', x)=|z_1+3k-1-a|+|z_2-b|=|z_1-c+3k|+|z_2-d+k-1|\geq 2k.\]   
%\note{Henry asks: If I want to try to prove Lemmas~\ref{lem:facets_3kor3k-1} and~\ref{lem:facets_3kor3k-1-1} at the same time, then is it correct to say that this paragraph changes based on the maximal simplices in the cyclic graph, and that the above paragraph doesn't change much and still follows from all the earlier lemmas in this section? If so, I may try this out.}
%\note{Ziqin says: Yes. I believe 5.5 and 5.6 follow from the previous lemmas and the clarifications of facets in the clique power of cyclic graphs.} 
%\note{John says: In the $n=3k-1$ case, I can see how, in the case that all elements in $\tau$ have the same first or second coordinate, then $\tau \in N_{3k-1,k}$. However, I'm having some difficulty in seeing how this would imply $\tau\in M_{3k-1,k}$ if $\tau$ contains $[x] =([a], [b])$ and $[y]=([c], [d])$ such that $[a]\neq [c]$ and $[b]\neq [d]$.  How would you apply Lemma 5.4 as $n-2k=k-1$? In this case, the inequalities $1\leq|a-c|\leq k-1$ and $1\leq|b-d|\leq k-1$ do not necessarily guarantee that $|a-c|<n-2k=k-1$ and $|b-d|<n-2k=k-1$. If the Lemma 5.4 conditions were instead $|x_1-y_1|\leq n-2k$ and $|x_2-y_2|\leq n-2k$, then I believe it would work.} \note{Henry says: John, good question; I am not sure.}

In the remaining case when all elements of $\tau$ have the same first coordinate or second coordinate, assume again that the first coordinate of all elements in $\tau$ is $[a]$.
The full subcomplex of $\vr{T_{3k-1,3k-1}}{k}$ whose vertices have first coordinate $[a]$ is isomorphic to the clique complex $\vr{C_{3k-1}}{k}$.
Of the maximal simplices in this copy of $\vr{C_{3k}}{k}$, only the ones in $N_{3k-1,k}$ are maximal in $\vr{T_{3k-1,3k-1}}{k}$.
\end{proof}

\end{document}
